\Article{233}
If a (primitive) form $F=(a,b,c)$ is such that two numbers $g$, $h$ can be found with $gg\equiv a$, $gh\equiv b$, $hh\equiv c$ for a given modulus $m$, we will call this form a quadratic residue of the number $m$, and $gx+hy$ a value of the expression $\surd(axx+2bxy+cyy)(\mathrm{mod.}\ m)$, or, more briefly, $(g,h)$ a value of the expression $\surd(a,b,c)$ or $\surd F(\mathrm{mod.}\ m)$. More generally, if a multiplier $M$, relatively prime to the modulus $m$, is such that we can have
\[
gg\equiv aM,\quad gh\equiv bM,\quad hh\equiv cM\quad (\mathrm{mod.}\ m)
\]
we will call $M\times(a,b,c)$ or $MF$ a quadratic residue of $m$, and $(g,h)$ a value of the expression $\surd M(a,b,c)$ or $\surd MF(\mathrm{mod.}\ m)$. For \textit{example}, the form $(3,1,54)$ is a quadratic residue of $23$, and $(7,10)$ is a value of the expression $\surd(3,1,54)(\mathrm{mod.}\ 23)$; similarly, $(2,-4)$ is a value of the expression $\surd5(10,3,17)(\mathrm{mod.}\ 23)$. We will show the use of these definitions below; here we note the following propositions:

\par
\SourceLeaf{242}
I. If $M(a,b,c)$ is a quadratic residue of the number $m$, this number divides the determinant of the form $(a,b,c)$. For if $(g,h)$ is a value of the expression $\surd M(a,b,c)$ (mod.\ $m$), or
\[
gg\equiv aM,\quad gh\equiv bM,\quad hh\equiv cM\ (\mathrm{mod.}\ m)
\]
then $bbMM-acMM\equiv0$, or $(bb-ac)MM$ is divisible by $m$. But since $M$ is assumed relatively prime to $m$, $bb-ac$ is also divisible by $m$.

II. If $M(a,b,c)$ is a quadratic residue of $m$, and $m$ is either a prime number or a power of a prime number, say $=p^\mu$, the individual character of the form $(a,b,c)$ with respect to the number $p$ is either $Rp$ or $Np$, according as $M$ is a residue or a nonresidue of $p$. This follows at once because both $aM$ and $cM$ are residues of $m$, or of $p$, and at least one of the numbers $a$, $c$ is not divisible by $p$ (Article 230).

Similarly, if (with the other conditions unchanged) $m=4$, either $1,4$ or $3,4$ is the individual character of the form $(a,b,c)$, according as $M\equiv1$ or $\equiv3$; a similar result holds for $m=8$ or a higher power of the number $2$. The individual character is $1,8;\ 3,8;\ 5,8;\ 7,8$ for the form $(a,b,c)$, according as $M\equiv1;\ 3;\ 5;\ 7\ (\mathrm{mod.}\ 8)$, respectively.

III. Conversely, if $m$ is a prime number or a power of an odd prime number $=p^\mu$, dividing the determinant $bb-ac$, and $M$ is either a residue or a nonresidue of $p$, according as the character of the form $(a,b,c)$ with respect to $p$ is $Rp$ or $Np$, respectively, then $M(a,b,c)$ is a quadratic residue of $m$. For when $a$ is not divisible by $p$, $aM$ is a residue of $p$ and therefore also of $m$; so if $g$ is a value of the expression $\surd aM$ (mod.\ $m$), and $h$ a value of the expression $\frac{bg}{a}$ (mod.\ $m$), we have $gg\equiv aM;\ ah\equiv bg$, and so
\[
agh\equiv bgg\equiv abM\quad\text{and}\quad gh\equiv bM
\]
and finally
\[
ahh\equiv bgh\equiv bbM\equiv bbM-(bb-ac)M\equiv acM
\]
and so $hh\equiv cM$, \textit{i.e.}, $(g,h)$ is a value of the expression $\surd M(a,b,c)$. But when $a$ is divisible by $m$, $c$ certainly is not; it is then easy to see that the same result follows if we take for $h$ a value of the expression $\surd cM$ (mod.\ $m$). For $g$ we take a value of the expression $\frac{bh}{c}$ (mod.\ $m$).

Similarly, we prove that if we have $m$ $=4$ and this number divides $bb-ac$, and we take the number $M$ either $\equiv1$ or $\equiv3$, according as $1,4$ or $3,4$ is the individual character
\SourceLeaf{243}
of the form $(a,b,c)$, then $M(a,b,c)$ is a quadratic residue of $m$. Likewise, if we have $m$ $=8$ or a higher power of $2$, by which $bb-ac$ is divisible, and we take $M$ to be $\equiv1;\ 3;\ 5;\ 7\ (\mathrm{mod.}\ 8)$, as the individual character of the form $(a,b,c)$ with respect to the number $8$ requires, then $M(a,b,c)$ is a quadratic residue of $m$.

IV. If the determinant of the form $(a,b,c)$ is given by $=D$, and $M(a,b,c)$ is a quadratic residue of $D$, all the individual characters of the form $(a,b,c)$, both with respect to each odd prime divisor of $D$ and with respect to the number $4$ or the number $8$ (if they divide $D$), can be read off at once from the number $M$. For \textit{example}, since $3(20,10,27)$ is a quadratic residue of $440$, namely $(150,9)$ is a value of the expression $\surd3(20,10,27)$ modulo $440$, and $3N5$, $3R11$, the characters of the form $(20,10,27)$ are $3,8;\ N5;\ R11$. Only the individual characters with respect to the numbers $4$ and $8$, whenever these do not divide the determinant, have no necessary connection with the number $M$

V. Conversely, if the number $M$, relatively prime to $D$, contains all the individual characters of the form $(a,b,c)$ (except the characters with respect to the numbers $4,8$, when these do not divide $D$), then $M(a,b,c)$ is a quadratic residue of $D$. For III shows that if $D$ is put in the form $\pm A^\alpha B^\beta D^\gamma\ldots$, where $A$, $B$, $C$, etc. are distinct prime numbers, then $M(a,b,c)$ is a quadratic residue of each of $A^\alpha$, $B^\beta$, $C^\gamma$, etc. So if a value of the expression $\surd M(a,b,c)$ modulo $A^\alpha$ is $(\mathfrak{A},\mathfrak{A}')$; modulo $B^\beta$, $(\mathfrak{B},\mathfrak{B}')$; modulo $C^\gamma$, $(\mathfrak{C},\mathfrak{C}')$, etc., and the numbers $g$, $h$ are determined so that $g\equiv\mathfrak{A},\mathfrak{B},\mathfrak{C}$, etc.; $h\equiv\mathfrak{A}',\mathfrak{B}',\mathfrak{C}'$, etc., for the moduli $A^\alpha$, $B^\beta$, $C^\gamma$, etc., respectively (Article 32), we readily see that $gg\equiv aM$, $gh\equiv bM$, $hh\equiv cM$ for all the moduli $A^\alpha$, $B^\beta$, $C^\gamma$, etc., and therefore also for the modulus $D$, which is their product.

VI. For these reasons, numbers such as $M$ will be called \textit{characteristic numbers} of the form $(a,b,c)$. By V, several such numbers can be found without difficulty as soon as all the individual characters of this form are known; the simplest ones, however, are usually found most easily by trial. Clearly, if $M$ is a characteristic number of a given primitive form with determinant $D$, all numbers congruent to $M$ modulo $D$ are characteristic numbers of the same form; forms in the same class, or even in different classes of the same genus, have the same characteristic numbers, so any characteristic
\SourceLeaf{244}
number of a given form can also be assigned to the whole class and genus; finally, $1$ is always a characteristic number of the principal form, class and genus, or every form of the principal genus is a residue of its own determinant.

VII. If $(g,h)$ is a value of the expression $\surd M(a,b,c)$ (mod.\ $m$), and $g'\equiv g$, $h'\equiv h$ (mod.\ $m$), then $(g',h')$ is also a value of the same expression. Such values can be regarded as \textit{equivalent}; but if $(g,h)$, $(g',h')$ are values of the same expression $\surd M(a,b,c)$, and we do not have both $g'\equiv g$, $h'\equiv h$ (mod.\ $m$), they must be counted as \textit{distinct}. Clearly, whenever $(g,h)$ is a value of such an expression, $(-g,-h)$ is one too, and it is easy to prove that these values are always distinct unless $m=2$. It is just as easy to prove that the expression $\surd M(a,b,c)$ (mod.\ $m$) cannot have more distinct values than two such (opposite) ones when $m$ is either an odd prime number, a power of an odd prime number, or $=4$; but when $m$ is $=8$ or a higher power of the number $2$, there are exactly four. From this, using VI, we easily deduce that if the determinant $D$ of the form $(a,b,c)$ is $=\pm2^\mu A^\alpha B^\beta\ldots$, where $A$, $B$, etc. denote distinct odd prime numbers whose number is given by $=n$, and $M$ is a characteristic number of this form, there are exactly $2^n$, $2^{n+1}$ or $2^{n+2}$ distinct values of the expression $\surd M(a,b,c)$ (mod.\ $D$), according as $\mu$ is $<2$, $=2$ or $>2$. For \textit{example}, there are sixteen values of the expression $\surd7(12,6,-17)$ (mod.\ $240$), namely $(\pm18,\mp11)$, $(\pm18,\pm29)$, $(\pm18,\mp91)$, $(\pm18,\pm109)$, $(\pm78,\pm19)$, $(\pm78,\pm59)$, $(\pm78,\mp61)$, $(\pm78,\mp101)$. We omit a fuller proof for brevity, since it is not particularly needed for what follows.

VIII. Finally, we observe that if two equivalent forms $(a,b,c)$, $(a',b',c')$ have determinant $D$ and characteristic number $M$, and the first passes into the second by the substitution $\alpha,\beta,\gamma,\delta$, then any value of the expression $\surd M(a,b,c)$, such as $(g,h)$, gives a value of the expression $\surd M(a',b',c')$, namely $(\alpha g+\gamma h,\beta g+\delta h)$. Anyone can work out the proof without difficulty.

\Needspace{12\baselineskip}
\begin{center}
\textit{Composition of forms.}
\end{center}

\Needspace{8\baselineskip}
\Article{234}
Having discussed the division of forms into classes, genera and orders, and explained the general properties that follow directly from these distinctions, we turn to another very important subject, which no one has yet
\SourceLeaf{245}
considered: the \textit{composition} of forms. At the start of this investigation, we insert the following lemma at once, so that we will not later have to interrupt the continuous sequence of proofs.

\textsc{Lemma.} \textit{Suppose we have four sequences of integers}
\[
\begin{gathered}a,a',a''\ldots a^n;\qquad b,b',b''\ldots b^n;\\c,c',c''\ldots c^n;\qquad d,d',d''\ldots d^n\end{gathered}
\]
\textit{with equally many terms (say $n+1$), such that}
\[
cd'-dc',\quad cd''-dc''\ \text{\textit{etc.}},\quad c'd''-d'c''\ \text{\textit{etc. etc.}}
\]
\textit{are respectively}
\[
=k(ab'-ba'),\quad k(ab''-ba'')\ \text{\textit{etc.}},\quad k(a'b''-b'a'')\ \text{\textit{etc. etc.}}
\]
\textit{or, in general,}
\[
c^\lambda d^\mu-d^\lambda c^\mu=k(a^\lambda b^\mu-b^\lambda a^\mu)
\]
\textit{where $k$ denotes a given integer; $\lambda$, $\mu$ are any two distinct integers between $0$ and $n$ inclusive, the larger being $\mu$}\HistoricalNote{*)}{We regard $a$ as $a^0$, $b$ as $b^0$, etc. --- The same equation clearly also holds when $\lambda=\mu$ or $\lambda>\mu$.}\textit{; suppose also that all the $a^\lambda b^\mu-b^\lambda a^\mu$ have no common divisor. Then four integers $\alpha,\beta,\gamma,\delta$ can be found such that}
\[
\begin{aligned}
\alpha a+\beta b&=c, &\alpha a'+\beta b'&=c', &\alpha a''+\beta b''&=c''\ \text{\textit{etc.}}\\
\gamma a+\delta b&=d, &\gamma a'+\delta b'&=d', &\gamma a''+\delta b''&=d''\ \text{\textit{etc.}}
\end{aligned}
\]
\textit{or, in general,}
\[
\alpha a^\nu+\beta b^\nu=c^\nu,\quad\gamma a^\nu+\delta b^\nu=a^\nu
\]
\textit{and then we have}
\[
\alpha\delta-\beta\gamma=k
\]

Since by hypothesis the numbers $ab'-ba'$, $ab''-ba''$, etc., $a'b''-b'a''$, etc. (of which there are $=\frac12(n+1)n$) have no common divisor, we can find as many other integers such that, when the former are multiplied by these respectively, the sum of the products is $=1$ (Article 40). Denote these multipliers by $(0,1)$, $(0,2)$, etc., $(1,2)$, etc., or, in general, denote the multiplier of $a^\lambda b^\mu-b^\lambda a^\mu$ by $(\lambda,\mu)$, so that
\[
\Sigma(\lambda,\mu)(a^\lambda b^\mu-b^\lambda a^\mu)=1
\]
(By the letter $\Sigma$ we mean the sum of all values of the expression that
\SourceLeaf{246}
follows it, obtained by giving $\lambda$, $\mu$ all distinct values between $0$ and $n$, with $\mu>\lambda$). Now if we set
\[
\begin{aligned}
\Sigma(\lambda,\mu)(c^\lambda b^\mu-b^\lambda c^\mu)&=\alpha, &\Sigma(\lambda,\mu)(a^\lambda c^\mu-c^\lambda a^\mu)&=\beta\\
\Sigma(\lambda,\mu)(d^\lambda b^\mu-b^\lambda d^\mu)&=\gamma, &\Sigma(\lambda,\mu)(a^\lambda d^\mu-d^\lambda a^\mu)&=\delta
\end{aligned}
\]
these $\alpha,\beta,\gamma.\delta$ have the required properties.

\textit{Proof.} I. If $\nu$ denotes any integer between $0$ and $n$, we have
\[
\begin{aligned}
\alpha a^\nu+\beta b^\nu
&=\Sigma(\lambda,\mu)(c^\lambda b^\mu a^\nu-b^\lambda c^\mu a^\nu+a^\lambda c^\mu b^\nu-c^\lambda a^\mu b^\nu)\\
&=\frac1k\Sigma(\lambda,\mu)(c^\lambda d^\mu c^\nu-d^\lambda c^\mu c^\nu)\\
&=\frac1k c^\nu\Sigma(\lambda,\mu)(c^\lambda d^\mu-d^\lambda c^\mu)\\
&=c^\nu\Sigma(\lambda,\mu)(a^\lambda b^\mu-b^\lambda a^\mu)=c^\nu
\end{aligned}
\]
And a similar calculation gives
\[
\gamma a^\nu+\delta b^\nu=d^\nu.\qquad\text{\textit{Q. E. P.}}
\]

II. Since we therefore have
\[
c^\lambda=\alpha a^\lambda+\beta b^\lambda,\quad c^\mu=\alpha a^\mu+\beta b^\mu
\]
we get
\[
c^\lambda b^\mu-b^\lambda c^\mu=\alpha(a^\lambda b^\mu-b^\lambda a^\mu)
\]
and similarly
\[
\begin{aligned}
a^\lambda c^\mu-c^\lambda a^\mu&=\beta(a^\lambda b^\mu-b^\lambda a^\mu)\\
d^\lambda b^\mu-b^\lambda d^\mu&=\gamma(a^\lambda b^\mu-b^\lambda a^\mu)\\
a^\lambda d^\mu-d^\lambda a^\mu&=\delta(a^\lambda b^\mu-b^\lambda a^\mu)
\end{aligned}
\]
From these formulas the values of $\alpha,\beta,\gamma,\delta$ can be found much more easily, provided we choose $\lambda$, $\mu$ so that $a^\lambda b^\mu-b^\lambda a^\mu$ is not $=0$. This is certainly possible, because by hypothesis all the $a^\lambda b^\mu-b^\lambda a^\mu$ have no common divisor, and so they cannot \textit{all} be $=0$. --- From the same equations, multiplying the first by the fourth and the second by the third, and subtracting, we obtain
\[
(\alpha\delta-\beta\gamma)(a^\lambda b^\mu-b^\lambda a^\mu)^2
=(a^\lambda b^\mu-b^\lambda a^\mu)(c^\lambda d^\mu-d^\lambda c^\mu)
=k(a^\lambda b^\mu-b^\lambda a^\mu)^2
\]
and so necessarily
\[
\alpha\delta-\beta\gamma=k.\qquad\text{\textit{Q. E. S.}}
\]
